\documentclass[11pt]{article}
\usepackage[a4paper,margin=25mm]{geometry}
\usepackage{amsmath,amssymb,amsthm,hyperref}
\usepackage{lmodern}
\newtheorem{theorem}{Theorem}
\newtheorem{lemma}{Lemma}
\newcommand{\TV}{\mathrm{TV}}
\title{An explicit recursive mixing bound for the adopted CYCLON kernel}
\author{Working research note}
\date{16 November 2026}
\begin{document}
\maketitle
\paragraph{Scope and status.}
This is a quantitative companion to \texttt{generic-proof.tex}, not a
replacement for its algorithm specification and elementary controller proofs.
The protocol is exactly the iid uniform-global-initiator, mandatory partner
removal, received-identity protection, minimum-age duplicate merge kernel.
Membership is fixed; $1\le l\le c\le n-1$; initial numerical ages are finite
nonnegative integers; the initial underlying undirected graph is connected.
Empty and underfull caches are permitted. Time counts individual global events.
The new recurrence and its strengthened counts have passed independent internal
adversarial review, conditional on the explicitly identified existing lemmas.
This is not formal verification or external peer review.

\paragraph{Important limitation.}
The resulting bound is explicit but extremely pessimistic. It is not a
polynomial bound, not an estimate of typical simulation behaviour, and not a
claim of near-optimal mixing. Its main improvement over the general qualitative
proof is to replace unspecified finite-family path maxima by computable
recurrences uniform over all initial ages. A useful rate remains open here.

\section{Inputs and notation}
Let $A(n,c,l)$ bound a controlled path to any prescribed labelled canonical
zero-age state $H(V,c)$, from every allowed initial state. A controller specifies
legal outcomes of the original kernel, not a changed algorithm. Define
$A(1,0,0)=0$. The root of $H$ has an empty cache and hence admits a no-op.
Let $F(n,c,l,s)$ bound first contact across a frozen cut whose active side has
$s$ vertices and caches at least one outside identity. Global connectivity is
not required for this cut statement.

The inputs from \texttt{generic-proof.tex} are:
\begin{itemize}
\item Length-one canonical controllers: \texttt{thm:generalone} and
\texttt{thm:mandatoryone}.
\item Length-two cut controller: \texttt{thm:firstcontact}.
\item Colored controllers: \texttt{lem:generalcolor} and
\texttt{lem:partialcolorsource}.
\item Reserved-anchor simulation: \texttt{lem:anchorreduction}.
\item Four-event merging gadget: \texttt{prop:anchorcanonical}.
\end{itemize}
We quantify their combinations below.
No implicit higher-length common-atom constants are used as inputs.

For $l=1$, set
\[
 A(n,1,1)=2n(n-1).
\]
For $c\ge2$, define $b_1=0$,
\[
 b_m=2m(m+c-1)+2b_{m-1},\qquad
 A(n,c,1)=n(n+c-1)+b_n.
\]
These are the explicit length-one controller horizons in the companion proof.
For $l=2$, the existing cut theorem gives
\[
 d=s(s+c),\quad j=(sc+1)(2s+c),\quad F(n,c,2,s)=d(j+1)+j.
\]
At length one we may use $F(n,c,1,s)=s+c$. Indeed, repeatedly initiate a
peer $p$ caching an outside identity. Until crossing, this identity cannot be
evicted: there are no extras. Each internal responder learns and retains $p$;
a repeat responder can reply with $p$, shrinking $C_p$. There are at most
$s-1$ first internal responders and at most $c$ shrinking replies before the
outside identity must be contacted.

\paragraph{Optional stronger base.}
The full manuscript \texttt{main.tex}, Theorem \texttt{thm:c2explicitmix}, proves
$A(n,2,2)=16n^2$ and minorization $(4n)^{-16n^2}$. This theorem's proof is not
contained in the extracted generic companion. One may either import this
explicitly identified extra input, or use the recurrences below also for
$c=l=2$. Both versions are valid; numerical illustrations below state which
version they use.

\section{Uniform source creation}
For $l\ge2$, put
\begin{align*}
 M(n,c,l)&=\max_{1\le m<n}
 A\bigl(m,\min(c-1,m-1),\min(l-1,m-1)\bigr),\\
 E(n,c,l)&=1+\max\left\{M(n,c,l),
 \max_{1\le s<n-1}F(n-1,c-1,l-1,s)\right\}.
\end{align*}
All non-singleton arguments are allowed parameters. The active-side maximum
is over nonempty proper subsets of the reduced ambient population.
For $l\ge3$, the structural cut construction then gives
\[
 F(n,c,l,s)=sE(n,c,l)\bigl((sc+1)c+1\bigr)+(sc+1)c.
\]
This is the existing minimum-holder raising and young-count routing proof;
its only inputs have strictly smaller shuffle length.

Define integer tenure bounds
\[
 t_0=1,\qquad t_j=1+j t_{j-1}\quad(1\le j\le c-1),\qquad
 K=(n-1)t_{c-1}.
\]
Thus $t_j=j!\sum_{i=0}^{j}1/i!$. Put
\begin{align*}
 D&=ncl^c,& R&=(nc+1)(nc^2+1),\\
 H&=R\bigl((K+1)(D+1)+1\bigr),& B&=(nc+1)H,\\
 S(n,c,l)&=(B+1)(n-1)E(n,c,l)+H.
\end{align*}
\begin{lemma}
For any specified $q$, a controlled path of at most $S(n,c,l)$ events,
never initiating $q$, makes $q$ incoming-free with $C_q\ne\varnothing$.
This holds for every $l\ge2$ and uniformly over all initial numerical ages.
\end{lemma}
\begin{proof}
Freeze $q$. At the current minimum stored $q$-age $k$, reserve $q$ in
requests and replies between minimum holders. The internal residual kernel
has capacity $c-1$ and length $l-1$. If a proposed residual oldest is younger
than $k$, contacting $q$ already removes a minimum copy. Otherwise use the
lower-length cut controller on an open residual component, or the
lower-length canonical controller on a closed component followed by its
root's owner contact. The cost of one phase is at most $E$.
Internal contacts preserve both minimum copies through $\min(k+1,k)=k$.
A terminal boundary or owner contact loses a minimum copy and creates none.
At most $n-1$ phases clear a level. Since no fresh $q$ is generated, its
minimum cannot decrease. After at most $B+1$ integer levels, either $q$ is
absent or every surviving copy has age greater than $B$.

Choose an empty one of the $nc+1$ age intervals $(jH,(j+1)H]$.
Mark ages below the interval young and those above it old. During $H$
events all young ages remain below all old ages. All $q$ copies are old.
The existing partial colored controller therefore applies without any
bound on the passive row of $q$.

Here is the sharpened count in that controller. On a joint plateau of
$F=\sum f_i$ and $V=\sum f_i^2$, every outside responder is old-free,
and the fixed young count strictly increases along a newly transmitted
$q$ copy. A tenure at a row with $f$ young entries has at most $c-f-1$
child-producing activations: $q$ occupies one old slot and each non-$q$
activation deletes a distinct old head, receiving only young entries.
The row cannot be passively updated during the tenure. This remains true
if the last such activation evicts $q$. Its descendants therefore number
at most $t_{c-f-1}$ including itself. There are at most $n-1$ initial
roots, giving $K$ tenures and hence at most $K$ owner contacts per plateau.
Between these contacts the original potential
$W=\sum_i o_i l^{c-f_i}$ strictly decreases and lies between $0$ and $D$.
At most $R$ plateaus occur. This proves the stated $H$ bound.

Erasing the final $q$ reference does not require its last loss to be an
owner contact. A passive row cannot become empty; if $q$ was initially
empty and never contacted, connectivity prevents disappearance of all
incoming references. Thus the final source row is nonempty.
\end{proof}

\section{Normalization recurrences}
For nonnegative integer $N$, let $Q(N;d,e)$ be the maximum of
\[
 \sum_{m\in\mathcal P}A\bigl(m,\min(d,m-1),\min(e,m-1)\bigr)
\]
over integer partitions $\mathcal P$ of $N$. Singleton terms are zero.
Let $Q_{\rm proper}$ restrict to partitions with at least two parts.
These are finite maxima over component sizes, not numerical states.

A universal zero anchor can be normalized in
\[
 U(n,c,l)=(c+3)M(n,c,l)+Q(n-1;c-1,l-1)+c+4(n-2)+1
\]
events. To see this, empty its owner using at most $c$ lower component
normalizations and $c$ contacts. Normalize initial residual components at
cost $Q$. Merge their roots with the four-event gadget at most $n-2$ times.
The same singleton root can be reused, so no intermediate normalization
is charged. One final residual normalization gives a canonical state.
Two further residual normalizations and one contact allow any prescribed
root and order. These account for the three extra $M$ terms.

Write $J(n,l)=A(n,n-1,l)$ for full capacity, retaining the length-one base.
For $l\ge2$, put
\[
 m=\left\lceil\frac{n-2}{l-1}\right\rceil,
\]
\[
 J(n,l)=(m+1)S(n,n-1,l)+m
       +J(n-1,\min(l,n-2))+1+U(n,n-1,l).
\]
At full capacity there are no evictions, and no underlying undirected edge
can disappear. Create a source $q$. If its complement has $k>1$ components,
$q$ caches a representative of every component. Its request to an oldest
$u$ can include representatives of $\min(l-1,k-1)$ other components,
creating that many complementary edges at the responder. There are enough
extra slots because $|C_q|\ge k$. Recreate a source and repeat; the new
edges cannot disappear. After at most $m$ merges and $m+1$ source creations,
the complement is connected. Normalize it with root $r$ different from
the frozen oldest $u$ of $q$. One contact $q\to u$ receiving $r^0$ creates
a universal zero anchor, since no eviction can remove $u$'s $r^0$. Apply $U$.
The case $n=2$ was already length one.

For proper capacity $2\le l\le c<n-1$, define
\begin{align*}
 A(n,c,l)={}&S(n,c,l)+Q_{\rm proper}(n-1;c,l)\\
 &+\max_{1\le s\le c}J(s+1,\min(l,s))
 +A(n-1,c,l)+2+U(n,c,l).
\end{align*}
This counts the source-extension proof directly. Normalize disconnected
complement components with source-neighbor labels highest. Normalize the
closed core consisting of $q$ and its cache using the independent
full-capacity bound. Normalize the now-connected complement, make the two
source contacts giving a universal zero, and apply $U$. If the complement
was already connected, the initial component and core steps can be skipped;
the bound still holds. Every part in $Q_{\rm proper}$ has size at most $n-2$.

\paragraph{Well-foundedness.}
First establish all smaller lengths. At a new length establish all $J(n,l)$
by population induction; its source and anchor inputs have smaller length,
and its remaining input has smaller population and full capacity. Then
establish proper-capacity $A(n,c,l)$ by population induction. There is no
same-parameter circularity. Optionally replacing the $c=l=2$ cases with the
stronger explicit manuscript base only shortens these controllers.

\section{Probability and sample bounds}
\begin{theorem}
Put
\[
 Z=nc\binom{c-1}{l-1}\binom cl
 \begin{cases}
 1,&c=n-1,\\
 \displaystyle\binom{l-1}{\lfloor(l-1)/2\rfloor}
 \binom l{\lfloor l/2\rfloor},&c<n-1.
 \end{cases}
\]
Let $T=A(n,c,l)$, $\gamma=Z^{-T}$, and
$\alpha=(n)_c\gamma$, where $(n)_c=n!/(n-c)!$. For every prescribed labelled
canonical atom $h$, uniformly over all allowed initial states,
\[
 P^T(x,\{h\})\ge\gamma,\qquad
 \sup_x\|P^t(x,\cdot)-\pi\|_{\TV}
 \le(1-\alpha)^{\lfloor t/T\rfloor}.
\]
Consequently for $0<\varepsilon<1$,
\[
 t_{\rm mix}(\varepsilon)
 \le T\left\lceil\frac{\log(1/\varepsilon)}{-\log(1-\alpha)}\right\rceil
 \le T\left\lceil\alpha^{-1}\log(1/\varepsilon)\right\rceil.
\]
\end{theorem}
\begin{proof}
Each prescribed event has conditional probability at least $p=Z^{-1}$.
The initiator contributes $1/n$ and oldest tie at least $1/c$. Request and
reply subset counts are at most $\binom{c-1}{l-1}$ and $\binom cl$:
when a row is shorter than its packet all choices are forced; otherwise
the relevant binomial is increasing in the row size. The active eligible
victim set has size at most $l-1$, and the passive one at most $l$.
Any fixed-cardinality subset count is bounded by the central binomial at
that maximal size. At full capacity no eviction occurs, so both victim
factors are one. This also implies the older bound $Z\le nc16^c$.
The controllers
are adaptive legal witnesses; multiplying conditional bounds is valid.
Pad their endpoint using its empty-root no-op, whose probability $1/n$
is at least $p$. This bound holds simultaneously for every distinct labelled canonical state.
In $H(V,c)$, the top $c$ vertices have distinct row sizes $0,\ldots,c-1$;
the other $n-c$ source vertices have identical full rows and no incoming
references. Its stabilizer therefore has size $(n-c)!$, so its relabelling
orbit has $(n)_c$ distinct states. Summing their individual bounds gives
\[
 P^T(x,\cdot)\ge\alpha\operatorname{Unif}(\operatorname{Orb}H).
\]
Here $T\ge n$, $Z\ge n$ and $(n)_c\le n^c$ with $c<n$, so
$0<\alpha<1$. This minorization gives the stated Doeblin
contraction, uniqueness and stationarity on the full numerical state space.
\end{proof}
At $l=1$, the proper-capacity formula gives $Z=nc^2$; at $c=1$ it gives
$Z=n$. At $c=l=2$ with $n>3$, $Z=4n$, recovering the sharper known
minorization when that extra horizon base is imported. At full shuffle
and full capacity all packets are forced, giving $Z=nc$. Other specialized
estimates can likewise be retained rather than replaced by a worse bound.

For an ordered local $b$-sample, let $a_b=\pi(|C_i|\ge b)>0$.
Stationary relabeling symmetry makes the sample conditional on availability
uniform, and conditioning yields
\[
 \|\mathcal L_x(S_i^{(b)}(t)\mid |C_i(t)|\ge b)
       -\operatorname{Unif}((V\setminus\{i\})_b)\|_{\TV}
 \le \frac{(1-\alpha)^{\lfloor t/T\rfloor}}{a_b}.
\]
The conditional law is considered when its conditioning event has positive
probability; the bound guarantees this once the full-state error is below
$a_b$. Exactly $n-b$ vertices of a canonical state have cache size at least $b$.
The uniform-orbit minorization therefore gives $a_b\ge\alpha(n-b)/n$.
It also gives this lower bound for actual availability at every $t\ge T$,
by applying the minorization to the final $T$ events.
For one identifier the stronger elementary bound
$a_1\ge\lceil(n-1)/c\rceil/n$ follows from connectivity, capacity and
stationary label symmetry. At full capacity the eventual monotone undirected
graph is complete, giving $a_1\ge(n-1)/n$. These connectivity-based bounds
refer to stationary availability, not an arbitrary transient fixed observer.
For the full-shuffle full-capacity case $l=c=n-1$, there is a stronger
all-request-size bound $a_b\ge(n-b)/n$, proved below.

\subsection*{Hole-DAG availability at full capacity and full shuffle}
Assume $l=c=n-1$, and define $D_p=V\setminus(C_p\cup\{p\})$ as an
identity set. If $p\to q$ is the actual contact, the full-packet, no-eviction
rule gives exactly
\[
 D_q'=D_p\cap D_q,\qquad D_p'=(D_p\cap D_q)\cup\{q\}.
\]
All other rows are unchanged. If the directed hole graph is acyclic before
this update, it stays acyclic. Its only new arc can be $p\to q$.
A new simple cycle would contain $p\to q\to r\to\cdots\to p$,
where the new $q\to r$ also belonged to the old row $D_p$.
Replacing its first two edges by the old $p\to r$ would give an old cycle.
There is no new two-cycle because $p\notin D_q'$. This is a contradiction.

The canonical state's hole graph is acyclic. Closure of this property and
convergence from that state show that the stationary law is supported on
hole DAGs. In a forward topological order $v_1,\ldots,v_n$, peer $v_j$ has
at most $n-j$ holes, so $|C_{v_j}|\ge j-1$. At least $n-b$ rows are
therefore available for request size $b$. Stationary relabelling symmetry
gives
\[
 a_b\ge\frac{n-b}{n}\qquad(1\le b\le n-1).
\]
The aggregate statement is deterministic on hole DAGs, but the fixed-observer
bound uses stationarity (or an exchangeable law). It is not asserted for
arbitrary initial labels at finite time. Thus any valid full-state bound
$d(t)$ gives conditional sample error at most $nd(t)/(n-b)$ in this case.
This improves conditioning, not the full-state mixing rate. The all-$n$
proof and a finite exact checker are recorded in
\texttt{docs/research/FULL-SHUFFLE-HOLE-DAG.md}.

\paragraph{Scale of the bound.}
With the stronger two-entry base, $(n,c,l)=(4,3,3)$ gives
\[T=54{,}446{,}588{,}124\]
even after the sharpened tenure and component counts.
The mixing bound then exponentiates $T$ through $\alpha^{-1}$.
This is a rigorous general upper bound but not a useful practical rate.
The gap from the existing $\Omega(n\log n)$ worst-case full-state lower
bound remains enormous. An improved probabilistic contraction argument,
rather than forcing one exact path, is needed to close it. Full-state,
conditional local-sample, and cross-time independence claims remain distinct.
\section{Sharp first-contact scale for a special capacity-two anchor}
\label{sec:c2-anchor-sharp}
This section concerns a first-contact time, not a mixing time. Assume
$n\ge3$ and $c=L=2$. Fix a peer $q$ whose cache is empty and suppose that
every other cache contains $q^0$. The possible remaining descriptor in
each such cache may have arbitrary nonnegative age. Let $\tau$ count global
events up to and including the first contact with $q$, and put $m=n-1$.
Define
\[
 b_n=n B(1/2,n-1)+\frac{n}{n-1},
\]
where $B$ is the beta integral. Then every such initial state satisfies
\[
 \mathbb E\tau\le b_n\le n\sqrt{\frac{\pi}{n-2}}+\frac{n}{n-1}.
\]
Moreover, if the outside residual descriptors initially form one directed
$m$-cycle, all with positive age, then for $n\ge4$
\[
 b_n-\frac{b_n^2}{n-3}\le\mathbb E\tau\le b_n.
\]
Consequently the supremum of $\mathbb E\tau$ over this anchored family is
$\sqrt{\pi n}+O(1)$ as $n\to\infty$.

\paragraph{Ready-count comparison.}
Until $\tau$, $q$ stays empty and its descriptor stays at age zero in all
other rows. Call a row ready when its residual descriptor is absent or has
age zero, and let $R$ be the number of ready outside rows. For an internal
contact $p\to r$, the responder becomes $\{q^0,p^0\}$; the initiator inherits
the responder's old residual descriptor, except that a self-reply is
removed. Thus $R$ never decreases. A nonready initiation raises $R$ by at
least one, and a ready initiation contacts $q$ with probability at least
$1/2$. The probabilities of growth and killing are therefore at least
$(m-R)/n$ and $R/(2n)$, respectively, and these events are disjoint.

Compare with a chain on $\{0,\ldots,m,\dagger\}$ which, from $r$, moves to
$r+1$ with probability $(m-r)/n$, is killed with probability $r/(2n)$, and
otherwise stays put. The absorbing state $\dagger$ is ordered above $m$.
Its kernel is stochastically increasing: for $r<s$, every nonabsorbing
output from $r$ is at most $s$, and its killing probability is smaller.
The actual ready-count kernel dominates this kernel, since both its
probability of reaching at least $r+1$ and its killing probability are
larger. Starting the comparison at zero therefore stochastically delays
killing. Write $T$ for its killing time.

\paragraph{Exact comparison mean.}
Poissonize global events at independent rate one. Realize the comparison
by $m$ independent rate-$1/n$ clocks and one inert clock for $q$. A vertex's
first tick makes it ready; each later tick kills with probability $1/2$.
Its probability of not killing by continuous time $u$ is
$2e^{-u/(2n)}-e^{-u/n}$. Independence gives survival
$(2e^{-u/(2n)}-e^{-u/n})^m$. Expected Poissonized time equals expected event
count, because the independent inter-event intervals have mean one. Hence
\[
 \mathbb ET=2n\int_0^1 y^{m-1}(2-y)^m\,dy
 =2n\int_0^1(1-v^2)^{m-1}(1+v)\,dv
 =nB(1/2,m)+\frac nm=b_n.
\]
The elementary inequality $(1-v^2)^{m-1}\le e^{-(m-1)v^2}$ gives the stated
upper bound. The identity
\[B(1/2,m)=\sqrt\pi\,\Gamma(m)/\Gamma(m+1/2)\]
and Stirling's formula give
$b_n=\sqrt{\pi n}+O(1)$.

\paragraph{Sharpness.}
Start the outside residuals in the positive-age directed $m$-cycle.
Before contact with $q$, residual identities obey capacity-one rotations:
$p\to r\to s$ becomes $r\to p\to s$, with a self-reply removed if $s=p$.
The unique directed cycle decreases in length by at most one per event.
An empty residual can first appear only when a two-cycle is canceled.
Starting at length $m$, this cannot happen in the first $m-2=n-3$ events.
Throughout those events all ready residuals are nonempty, so their contact
probability is exactly $1/2$; a nonready initiation raises $R$ exactly by
one, and a surviving ready initiation preserves $R$. Thus the killed
ready-count process has exactly the comparison law through $h=n-3$ events.
The comparison's monotonicity bounds its expected remaining lifetime from
any transient state by $b_n$. For $n\ge4$, Markov's inequality therefore
implies
\[
 \mathbb E(T-h)_+\le\Pr(T>h)b_n\le b_n^2/h.
\]
Equality of survival through $h$ now yields
$\mathbb E\tau\ge\mathbb E\min(T,h)\ge b_n-b_n^2/(n-3)$.
The last subtraction is $O(1)$, proving the claim. Nothing in this argument
asserts a general-cut contact bound, a mixing bound, or an extension to
capacity three.

\end{document}
